\documentclass[11pt,a4paper]{article}

\usepackage{amsmath,amssymb,amsthm}
\usepackage{hyperref}

\newtheorem{theorem}{Theorem}[section]
\newtheorem{proposition}[theorem]{Proposition}

\DeclareMathOperator{\Gal}{Gal}

\title{Supplementary Material:\\
  Computational Galois Cohomology}
\author{}
\date{\today}

\begin{document}

\maketitle

\section{Smith Normal Form Algorithm}\label{sec:snf}

We implement a standard Smith normal form algorithm over $\mathbb{Z}$
with the following modifications for numerical stability:
\begin{enumerate}
  \item Pivot selection: choose the entry with smallest absolute value.
  \item Divisibility repair: if $d_i \nmid d_{i+1}$ after initial
    elimination, perform column addition and re-eliminate.
\end{enumerate}

Complexity: $O(n^3 \log M)$ where $M$ bounds the matrix entries.

\section{Coboundary Matrix Dimensions}\label{sec:dimensions}

For $|G| = m$ and $\text{rank}(M) = r$:
\begin{itemize}
  \item $\dim C^n(G, M) = r \cdot m^n$
  \item $d^n$ is an $(r m^{n+1}) \times (r m^n)$ matrix
  \item Total nonzero entries in $d^n$: at most $(n+2) \cdot r \cdot m^{n+1}$
\end{itemize}

For $|G| = 4$ (Klein four), $r = 1$, $n = 2$: $d^2$ is $64 \times 16$.

\section{Detailed Computation: $H^2(C_2 \times C_2, \mathbb{Z})$}

The Klein four group $V_4 = \{e, a, b, ab\}$ with trivial action on $\mathbb{Z}$.

$C^1(V_4, \mathbb{Z}) = \mathbb{Z}^4$ (one coordinate per group element).
$C^2(V_4, \mathbb{Z}) = \mathbb{Z}^{16}$ (indexed by pairs from $V_4$).

The $d^1$ matrix is $16 \times 4$. Its kernel (2-cocycles) has rank 10.
The image of $d^0$ (coboundaries in $C^1$) has rank 3.
The image of $d^1$ (coboundaries in $C^2$) has rank 8 among the 10 cocycles.

Result: $H^2 \cong \mathbb{Z}^{10}/\text{(rank-8 subgroup)}$ with
Smith invariants $[1,1,1,1,1,1,1,1,2,2]$ giving
$(\mathbb{Z}/2\mathbb{Z})^2$.

\section{Lean 4 Proof Strategy}\label{sec:lean}

Our Lean proofs reference Mathlib's \texttt{groupCohomology} API:
\begin{itemize}
  \item \texttt{RepresentationTheory.GroupCohomology.Basic} for $H^n$ definition
  \item \texttt{RepresentationTheory.GroupCohomology.Resolution} for bar resolution
  \item \texttt{RepresentationTheory.GroupCohomology.LowDegree} for $H^0$, $H^1$ characterizations
\end{itemize}

The $d^2 = 0$ property follows from the general theory of derived functors
(specifically, that the bar resolution gives a chain complex).

\end{document}
